% TouchModule
\section{【中級】TouchModule — バスオブジェクト共有パターン}

\begin{patternbox}{TouchModule で学べるパターン}
\begin{itemize}
    \item 外部ライブラリのドライバオブジェクト共有（TFT\_eSPIをTftModuleと共有）
    \item 入出力分離の原則（タッチ入力とLCD出力を別モジュールに分ける）
    \item 空のConfig構造体（設定がplatformio.iniで管理される場合）
\end{itemize}
\end{patternbox}

\subsection{ドライバオブジェクトの共有}

\begin{lstlisting}[style=cppstyle, caption={TouchModule.h — TFT\_eSPIの前方宣言}]
class TFT_eSPI;  // TftModuleと共有するドライバ

class TouchModule : public IModule {
private:
    TouchConfig _config;
    TFT_eSPI*   _tft;  // コンストラクタで受け取る共有ドライバ
public:
    TouchModule(const TouchConfig& config, TFT_eSPI* tft);
    // ...
};
\end{lstlisting}

\begin{lstlisting}[style=cppstyle, caption={main.cpp — ドライバの共有}]
static TFT_eSPI tftDriver;  // 1つのインスタンスを生成

// 入力と出力の両モジュールが同じドライバを共有
TouchModule touchModule(TOUCH_CONFIG, &tftDriver);
TftModule   tftModule(TFT_CONFIG, &tftDriver);
\end{lstlisting}

\begin{pointbox}[通信バスパターンの応用]
I2C/SPIのバスインスタンスだけでなく、外部ライブラリのドライバオブジェクトも同じパターンで共有できる。
TFT\_eSPIはSPIバスの抽象化であり、タッチ機能もSPIを共有するため、
1つのドライバインスタンスをTouchModuleとTftModuleの両方に渡す。
\end{pointbox}

\begin{pointbox}[入出力分離の実例]
TFTディスプレイはタッチ入力とLCD出力の両方を持つが、
\textbf{TouchModule（入力）とTftModule（出力）に分離}している。
入力フェーズでタッチ座標を読み取り、出力フェーズで画面を描画する。
これにより各モジュールの責務が「入力のみ」「出力のみ」に限定される。
\end{pointbox}
